\documentclass[11pt]{article}

\usepackage[margin=1in]{geometry}
\usepackage[T1]{fontenc}
\usepackage[utf8]{inputenc}
\usepackage{lmodern}
\usepackage{microtype}
\usepackage{amsmath,amssymb,amsthm}
\usepackage{booktabs}
\usepackage{tabularx}
\usepackage{enumitem}
\usepackage{xcolor}
\usepackage{hyperref}

\hypersetup{
  colorlinks=true,
  linkcolor=blue!55!black,
  citecolor=blue!55!black,
  urlcolor=blue!55!black,
  pdftitle={Coase--Information Theory},
  pdfauthor={Brian Guarraci}
}

\setlist[itemize]{topsep=4pt,itemsep=2pt,leftmargin=1.5em}
\setlist[enumerate]{topsep=4pt,itemsep=2pt,leftmargin=1.7em}
\renewcommand{\arraystretch}{1.15}

\newtheorem{proposition}{Proposition}
\newtheorem{definition}{Definition}
\newtheorem{remark}{Remark}

\newcommand{\E}{\mathbb{E}}
\newcommand{\calB}{\mathcal{B}}
\newcommand{\calX}{\mathcal{X}}
\newcommand{\calA}{\mathcal{A}}
\newcommand{\calM}{\mathcal{M}}
\newcommand{\calO}{\mathcal{O}}
\newcommand{\calT}{\mathcal{T}}

\title{\textbf{Coase--Information Theory}\\
\large Agent Boundaries, Organizational Agility, and Software Agents}
\author{Brian Guarraci}
\date{Working paper draft \\ June 2026}

\begin{document}
\maketitle

\begin{abstract}
This paper proposes a Coase--Information Theory of agent boundaries.
Classical transaction-cost economics explains the firm as an alternative to
market coordination when internal organization is cheaper than contracting. I
generalize this account by treating an agent as any bounded
information-processing unit: an individual, team, organization, firm, vendor,
software agent, or coalition of these. Boundaries form when agents coalesce into
shared representations and control loops; boundaries dissolve or split when the
cost of shared coordination exceeds its reduction in uncertainty, delay, error,
or surprise. The paper decomposes transaction and coordination costs into
information-processing terms: sensing, communication, interpretation,
monitoring, delay, error, and misalignment. It defines organizational agility as
the rate at which decision-relevant information about the world becomes
effective action. Software agents matter because they alter the cost surface of
coordination: they can reduce communication, interpretation, and execution
latency for structured tasks while introducing new governance and misalignment
costs. This creates an agent-mediated design region between hierarchy and
market, and a broader theory of when agents should merge, remain separate, or
coordinate through protocols.
\end{abstract}

\section{Introduction}

Coase framed the firm as a response to the costs of using the price mechanism:
firms exist because contracting, searching, negotiating, monitoring, and
enforcing every transaction through the market can be more costly than directing
work internally \cite{coase1937}. Modern organizations increasingly coordinate
through software systems, APIs, observability pipelines, ticket queues, chat
systems, workflow engines, and software agents. These tools do not merely lower
generic transaction costs. They change how information moves, how decisions are
formed, and how quickly actions can be coordinated.

This suggests a sharper question:

\begin{quote}
Given a changing environment, what agent architecture best converts
distributed, noisy information into coordinated action under latency, cost, and
governance constraints?
\end{quote}

The answer determines not only the boundary of the firm but also the boundaries
of teams, vendors, platforms, workflows, and software agents. It determines
which agents should coalesce into a larger agent, which should split apart, and
which should remain separate but coupled through protocols.

The central claim is that the firm is one instance of a broader
\emph{agent-boundary problem}. A person, team, department, firm, vendor network,
or software agent can each be treated as an agent when it has observations,
internal state, policies, actions, and a boundary. Organizational structure is a
communication graph. Roles, dashboards, runbooks, APIs, metrics, and managers
are compression mechanisms. Software agents are organizational nodes with
observation channels, message protocols, policies, tool interfaces, and
governance envelopes. The economic performance of an organization depends on
the fidelity, speed, and cost with which latent world states are converted into
action.

The contribution is not the broad claim that firms process information. That
intuition has deep roots in Hayek, Simon, Galbraith, and distributed decision
theory \cite{hayek1945,simon1947,galbraith1973,marschakradner1972}. The
contribution is a formal decomposition of boundary formation, organizational
cost, and agility that can be operationalized in software and workflow traces.

\section{Formal Model}

Let the external or internal environment generate a latent state
\[
X_t \in \calX
\]
at time \(t\). The state may include customer demand, a production incident, a
supply shock, a workflow backlog, a vendor outage, or any other
decision-relevant facts.

An organization is a directed graph
\[
G=(V,E),
\]
where nodes \(i \in V\) are humans, teams, services, vendors, software agents,
or composite groups. Edges \((i,j)\in E\) are communication or coordination
channels.

The word \emph{agent} is used recursively. A node may be a primitive actor, such
as a person or software service, or a composite actor, such as a team, firm,
platform, vendor network, or incident-response group. Let \(V_0\) be the set of
primitive actors and let
\[
\calB=\{B_1,\ldots,B_m\}
\]
be a partition of \(V_0\) into bounded agents. Each block \(B_\ell\) is an agent
when its members share enough representation, policy, governance, and action
interfaces to be treated as one decision-making unit at the relevant level of
analysis.

Each node receives a local observation
\[
Y_t^i \sim P_i(Y\mid X_t)
\]
and may send messages
\[
M_t^{ij}\in\calM_{ij}
\]
to neighboring nodes. Messages may be unstructured human communication,
structured tickets, API calls, database writes, workflow events, or
agent-generated summaries.

Each node maintains a compressed local representation
\[
Z_t^i=f_i(Y_{\leq t}^i,M_{\leq t}^{*i}),
\]
where \(M_{\leq t}^{*i}\) denotes messages received by node \(i\). Node \(i\)
chooses an action according to a local policy
\[
A_t^i \sim \pi_i(Z_t^i).
\]
The organization produces a joint action
\[
A_t=g(A_t^1,\ldots,A_t^n)
\]
and receives utility \(U(X_t,A_t)\).

The general organizational design problem is
\[
\max_{\calB,G,\Pi,F}
\E\left[\sum_t U(X_t,A_t)\right]
-C_{\mathrm{comm}}
-C_{\mathrm{comp}}
-C_{\mathrm{delay}}
-C_{\mathrm{error}}
-C_{\mathrm{misalign}}
-C_{\mathrm{monitor}},
\]
where \(\Pi=\{\pi_i\}\) is the set of local policies and \(F=\{f_i\}\) is the
set of compression or representation maps.

\begin{table}[t]
\centering
\caption{Information-processing decomposition of coordination cost}
\begin{tabularx}{\textwidth}{@{}lX@{}}
\toprule
Term & Interpretation \\
\midrule
\(C_{\mathrm{comm}}\) & Meetings, messages, handoffs, routing, API calls, protocol overhead. \\
\(C_{\mathrm{comp}}\) & Human or machine effort to analyze, summarize, decide, or execute. \\
\(C_{\mathrm{delay}}\) & Loss from acting after the state has changed or the opportunity has decayed. \\
\(C_{\mathrm{error}}\) & Loss from stale, incomplete, distorted, or misrouted information. \\
\(C_{\mathrm{misalign}}\) & Loss from local actions that do not compose into good global action. \\
\(C_{\mathrm{monitor}}\) & Governance, audit, review, compliance, and oversight burden. \\
\bottomrule
\end{tabularx}
\end{table}

This decomposition is the Coase--Information move. Transaction cost is not a
primitive. It is partly an information-processing quantity.

\section{Endogenous Agent Boundaries}

The boundary of an agent determines which observations, representations, and
actions are internal to a shared control loop and which must cross an interface.
Internal coordination can preserve richer context and reduce surprise, but it
also creates communication, governance, and computation costs. External
coordination can reduce internal burden and preserve modularity, but it can
increase interpretation, monitoring, contracting, and delay costs.

Let \(Z_t^B\) denote the internal representation held by a bounded agent \(B\).
One idealized measure of surprise is
\[
S_t^B=-\log P_B(X_t\mid Z_t^B),
\]
the degree to which the true state is unexpected under the agent's
representation. In decision settings, the same role can be played by prediction
error or action loss \(L(X_t,A_t^B)\). The boundary-design problem can be stated
as choosing a partition \(\calB\) that trades off expected surprise and action
loss against coordination cost:
\[
\min_{\calB,G,F,\Pi}
\E\left[\sum_t S_t(\calB)+L(X_t,A_t)\right]+C(\calB,G,F,\Pi).
\]

Two agents should coalesce when the merged boundary lowers expected surprise or
action loss more than it raises internal coordination and governance cost:
\[
\E[S_{\mathrm{merged}}+L_{\mathrm{merged}}]+C_{\mathrm{merged}}
<
\E[S_{\mathrm{separate}}+L_{\mathrm{separate}}]+C_{\mathrm{separate}}.
\]
They should split when the inequality reverses. The boundary of the firm is a
special case of a more general agent-boundary problem: where should the system
draw the line between shared representation and protocol-mediated coordination?

\begin{definition}[Agent boundary]
An agent boundary is a partition of actors, observations, representations,
policies, and actions into an internal control loop plus external interfaces.
The boundary is efficient when moving it would not lower expected surprise,
action loss, and coordination cost.
\end{definition}

\section{Organizational Agility}

Speed alone is not agility. An organization can act quickly and incorrectly.
Accuracy alone is not agility either. An organization can make a correct
decision too late.

Let \(\tau_t\) be the end-to-end latency from signal availability to coordinated
action. Let decision quality be proxied by the mutual information between the
latent world state and organizational action:
\[
Q_t=I(X_t;A_t).
\]
Define organizational agility as
\[
\calA_t=\frac{I(X_t;A_t)}{\tau_t}.
\]
A cost-adjusted version is
\[
\calA_t^{(c)}=\frac{I(X_t;A_t)}{\tau_t C_t},
\]
where \(C_t\) is the resource cost of producing the action.

This definition captures the core intuition: higher information fidelity
improves agility; lower latency improves agility; lower cost improves
cost-adjusted agility; and delay and distortion can both destroy organizational
performance.

\subsection{Practical Proxies}

In most organizational traces, \(I(X_t;A_t)\) will not be directly observable.
The empirical version should use proxies, as shown in Table~\ref{tab:proxies}.

\begin{table}[t]
\centering
\caption{Operational proxies for formal quantities}
\label{tab:proxies}
\begin{tabularx}{\textwidth}{@{}lX@{}}
\toprule
Formal quantity & Operational proxy \\
\midrule
\(I(X_t;A_t)\) & Action correctness, SLA success, no rework, severity reduction, expert quality score. \\
\(S_t\) or prediction surprise & Unexpected escalation, reopen, SLA miss, anomaly, decision reversal, incident severity surprise. \\
\(\tau_t\) & Time to triage, time to decision, time to resolution. \\
\(C_t\) & Handoffs, people involved, messages, meetings, queue time, tool calls. \\
\(C_{\mathrm{error}}\) & Reopen rate, rollback rate, bad routing, decision reversal. \\
\(C_{\mathrm{monitor}}\) & Approvals, audits, review steps, escalation requirements. \\
\bottomrule
\end{tabularx}
\end{table}

The theory should be honest: mutual information is the idealized concept;
workflow telemetry provides measurable proxies.

\section{Compression and Organizational Structure}

Organizations rarely transmit the full state \(X_t\). Instead, they construct
internal representations
\[
X_t \rightarrow Z_t \rightarrow A_t.
\]
Roles, dashboards, metrics, status reports, tickets, runbooks, APIs, schemas,
plans, OKRs, approvals, and incident summaries are compression mechanisms. They
reduce communication burden, but they can also destroy decision-relevant
information.

The organizational design problem has the structure of a rate-distortion
problem:
\[
\min_f C_{\mathrm{repr}}(f)
+\lambda\,\E[L(X_t,A_t(f(X_t)))],
\]
where \(C_{\mathrm{repr}}(f)\) is the cost of maintaining and communicating
representation \(f\), and \(L\) is decision loss. A dashboard that compresses a
complex operational state into one green/yellow/red indicator is valuable if it
preserves the information needed for action. It is dangerous if it hides the
structure needed for diagnosis.

\section{Software Agents as Organizational Nodes}

A software agent can be modeled as a node
\[
\mathrm{Agent}_i=(\calO_i,\calM_i,\pi_i,\calT_i,\Gamma_i),
\]
where \(\calO_i\) is the set of observation channels, \(\calM_i\) is the
message schema or protocol, \(\pi_i\) is the local decision policy, \(\calT_i\)
is the set of tools or actuators, and \(\Gamma_i\) is the governance envelope:
permissions, constraints, auditability, and escalation rules.

Agents change the feasible set of organizational designs. They can observe
continuously, summarize and route information, enforce schemas, call tools,
monitor queues, escalate uncertainty, and execute routine actions. But agentic
coordination is not automatically superior. It may increase local-global
misalignment, proxy optimization, monitoring burden, structured propagation of
bad assumptions, and compression loss from over-summary.

The relevant comparison is conditional:

\begin{quote}
Agents improve organizational agility when tasks are sufficiently structured,
protocols are sufficiently expressive, and governance is strong enough to
control misalignment.
\end{quote}

\section{Boundary of the Firm}

The boundary of the firm is one level of the broader boundary-selection problem.
At any level, the question is whether coordination should happen through a
shared internal representation or through an interface between separate agents.

Let \(k\) index a coordination form:
\[
k\in\{\mathrm{internal},\mathrm{market},\mathrm{hybrid}\}.
\]
Each form has cost
\[
C_k=C_{\mathrm{search},k}
+C_{\mathrm{contract},k}
+C_{\mathrm{comm},k}
+C_{\mathrm{interp},k}
+C_{\mathrm{monitor},k}
+C_{\mathrm{delay},k}
+C_{\mathrm{error},k}
+C_{\mathrm{misalign},k}.
\]
The efficient boundary is selected by
\[
k^*=\arg\min_k C_k-\E[U_k].
\]

Classical hierarchy reduces some contracting and search costs but can increase
internal communication delay and bureaucratic compression. Market coordination
can reduce internal burden but increase search, contracting, monitoring, and
interpretation costs. Agent-mediated hybrid coordination can reduce
cross-boundary interpretation and delay if protocols and governance are strong
enough.

The claim is not that agents make firms smaller or markets always better. The
claim is that software agents and machine-readable protocols alter enough cost
terms to create a new efficient region between classic hierarchy and classic
market contracting. The deeper claim is that firms, teams, vendors, and
software agents are all boundary choices over information-processing units.
Coalescence is valuable when shared state reduces surprise, delay, and
misalignment more than it raises internal coordination cost. Splitting is
valuable when modularity and local autonomy reduce cost more than they increase
interface surprise.

\section{Propositions}

\begin{proposition}[Protocol informativeness]
For a fixed organizational graph and decision policy class, improving protocol
quality weakly increases achievable expected utility when it makes received
representations more informative about the latent state without increasing
delay or cost.
\end{proposition}

\noindent\textit{Intuition.} A machine-readable schema, API, or structured
ticket that preserves relevant state is a less garbled signal. Better
downstream decisions become feasible.

\begin{proposition}[Latency under volatility]
Suppose the relevance of information decays with environmental hazard rate
\(\lambda\). If a decision is taken after latency \(\tau\), the value of
otherwise correct information is approximately discounted by
\(e^{-\lambda \tau}\). Therefore, the marginal value of latency reduction rises
with volatility.
\end{proposition}

\noindent\textit{Sketch.} If the probability that the relevant state remains
stable after delay \(\tau\) is \(e^{-\lambda \tau}\), then effective information
available for action is proportional to \(I_0 e^{-\lambda \tau}\). Higher
\(\lambda\) makes delay more costly.

\begin{proposition}[Agentic routing]
For structured or semi-structured workflows, agent nodes can increase
organizational agility when their reduction in routing, interpretation, and
execution latency exceeds their added monitoring and misalignment costs.
\end{proposition}

\noindent A sufficient condition is
\[
\Delta C_{\mathrm{delay}}
+\Delta C_{\mathrm{comm}}
+\Delta C_{\mathrm{interp}}
>
\Delta C_{\mathrm{monitor}}
+\Delta C_{\mathrm{misalign}}
+\Delta C_{\mathrm{error}}.
\]

\begin{proposition}[Boundary coalescence and splitting]
When cross-boundary communication, interpretation, and monitoring become
programmable, the efficient boundary shifts toward the partition of agents that
minimizes expected surprise, action loss, and coordination cost.
\end{proposition}

\noindent\textit{Intuition.} APIs, schemas, shared observability, and agents
reduce the information costs that traditionally made external coordination
expensive. This can make splitting more attractive for modular work, while high
interdependence can still make coalescence into a larger agent more efficient.

\begin{proposition}[Governance limit]
Agentic decentralization has diminishing or negative returns when governance
quality is low because autonomy can increase misalignment and monitoring costs
faster than it reduces delay.
\end{proposition}

\noindent\textit{Intuition.} A fast local optimizer is valuable only when its
actions compose into the global objective.

\section{Executable Illustrations}

The companion repository includes five static web applications. The apps are
executable diagrams, not calibrated causal estimates.

\begin{itemize}
  \item \textbf{Topology Lab} compares hierarchy, market coordination, and
  agentic mesh under shared environmental assumptions.
  \item \textbf{Incident Room} simulates deployment incidents, customer
  escalations, lead routing failures, and vendor outages.
  \item \textbf{Boundary Explorer} compares in-house, outsourced, and hybrid
  agentic boundary choices using decomposed cost terms.
  \item \textbf{Support Router} shows when structured intake, policy coverage,
  agent autonomy, and knowledge quality improve first-touch accuracy, handoffs,
  and resolution time.
  \item \textbf{Workflow Trace Mapper} accepts a CSV event trace and computes
  triage latency, resolution time, handoffs, agent-assisted share, rework, and
  an agility proxy.
\end{itemize}

\section{Empirical Agenda}

A first empirical study should use incident response or support routing traces.
The minimal research design is:
\begin{enumerate}
  \item collect workflow traces,
  \item map events to cases,
  \item compute latency, handoffs, surprise, and action-quality proxies,
  \item compare human-heavy, protocolized, and agent-assisted workflows,
  \item estimate whether agentic coordination improves agility after
  controlling for case complexity.
\end{enumerate}

A practical agility proxy is
\[
\widehat{\calA}
=\frac{Q}{1+T_{\mathrm{resolution}}}\cdot\frac{1}{1+\alpha H},
\]
where \(Q\) is quality, \(T_{\mathrm{resolution}}\) is resolution time, \(H\) is
handoff count, and \(\alpha\) is a handoff penalty. This proxy is not a direct
mutual-information estimate. It is a practical measurement bridge.

To study coalescing and splitting, compare whether a workflow performs better
when work is handled inside one bounded agent, split across separate agents
with an explicit protocol, or coordinated through a hybrid boundary such as an
agent-managed vendor or escalation path. The empirical question is whether the
boundary change lowers surprise and action loss enough to justify any added
coordination or monitoring cost.

\section{Limitations}

This framework has several limitations.

First, mutual information is not directly observable in most organizations. The
paper must distinguish the idealized formal quantity from operational proxies.

Second, agents do not universally improve coordination. They help most when
workflows are structured, observability is high, and governance is clear.

Third, some organizational knowledge is tacit and difficult to encode into
protocols. Over-protocolization can compress away context.

Fourth, the current apps are toy models. They illustrate mechanisms but do not
yet estimate causal effects.

Fifth, organizational utility is multi-objective. Speed, accuracy, compliance,
safety, morale, and strategic learning may trade off.

\section{Conclusion}

Coase--Information Theory reframes the firm as an information architecture and
then generalizes that idea. The key economic question is not only whether work
should be coordinated by markets or hierarchy, but which agents should
coalesce, which should split, and which should coordinate through protocols so
that distributed observations become coordinated action.

Organizational agility is the rate of this transformation. Software agents
matter because they can change the latency, fidelity, and cost of coordination.
Their importance is not simply that they automate tasks. It is that they alter
where boundaries should be drawn: inside firms, across teams, between vendors,
and between human and machine actors.

The next generation of AI-native companies will be designed not merely around
people using AI tools, but around software-mediated sensing, routing, decision,
and execution loops. The purpose of Coase--Information Theory is to provide a
language for designing and measuring those loops.

\begin{thebibliography}{9}

\bibitem{coase1937}
Coase, R. H. (1937).
\newblock The nature of the firm.
\newblock \emph{Economica}, 4(16), 386--405.

\bibitem{shannon1948}
Shannon, C. E. (1948).
\newblock A mathematical theory of communication.
\newblock \emph{Bell System Technical Journal}, 27(3), 379--423.

\bibitem{hayek1945}
Hayek, F. A. (1945).
\newblock The use of knowledge in society.
\newblock \emph{American Economic Review}, 35(4), 519--530.

\bibitem{simon1947}
Simon, H. A. (1947).
\newblock \emph{Administrative Behavior}.
\newblock Macmillan.

\bibitem{marschakradner1972}
Marschak, J., and Radner, R. (1972).
\newblock \emph{Economic Theory of Teams}.
\newblock Yale University Press.

\bibitem{galbraith1973}
Galbraith, J. R. (1973).
\newblock \emph{Designing Complex Organizations}.
\newblock Addison-Wesley.

\bibitem{williamson1985}
Williamson, O. E. (1985).
\newblock \emph{The Economic Institutions of Capitalism}.
\newblock Free Press.

\end{thebibliography}

\end{document}
